\chapter{Method}\label{ch:method}

\section{System Overview}
Describe the overall design. A block diagram helps (\cref{fig:blocks}).

\begin{figure}[ht]
  \centering
  \begin{tikzpicture}[node distance=1.6cm, every node/.style={font=\small},
      blk/.style={draw=atl-primary, fill=atl-light, rounded corners=2pt, minimum width=2.4cm, minimum height=0.9cm, thick},
      arr/.style={-{Stealth[length=2.2mm]}, thick, atl-secondary}]
    \node[blk] (a) {Input};
    \node[blk, right=of a] (b) {Processing};
    \node[blk, right=of b] (c) {Output};
    \draw[arr] (a) -- (b); \draw[arr] (b) -- (c);
  \end{tikzpicture}
  \caption{Block diagram of the system.}\label{fig:blocks}
\end{figure}

\section{Implementation}
Show the key algorithm or code.
\begin{codebox}[title=algorithm.py]{Python}
def moving_average(x, n):
    return [sum(x[i:i + n]) / n for i in range(len(x) - n + 1)]
\end{codebox}

\section{Chapter Summary}
Summarise the design decisions.
